
\documentclass{article}
\usepackage{colortbl}
\usepackage{makecell}
\usepackage{multirow}
\usepackage{supertabular}

\begin{document}

\newcounter{utterance}

\twocolumn

{ \footnotesize  \setcounter{utterance}{1}
\setlength{\tabcolsep}{0pt}
\begin{supertabular}{c@{$\;$}|p{.15\linewidth}@{}p{.15\linewidth}p{.15\linewidth}p{.15\linewidth}p{.15\linewidth}p{.15\linewidth}}

    \# & $\;$A & \multicolumn{4}{c}{Game Master} & $\;\:$B\\
    \hline 

    \theutterance \stepcounter{utterance}  

    & & \multicolumn{4}{p{0.6\linewidth}}{\cellcolor[rgb]{0.9,0.9,0.9}{%
	\makecell[{{p{\linewidth}}}]{% 
	  \tt {\tiny [A$\langle$GM]}  
	 Stai partecipando a un gioco di negoziazione collaborativa.\\ \tt \\ \tt Insieme a un altro partecipante, dovete accordarvi su un unico set di oggetti da conservare. Ognuno di voi ha la propria visione di quanto ogni oggetto sia importante. Non sai come l'altro partecipante valuta gli oggetti. Inoltre, ti viene fornito lo sforzo che ogni oggetto richiede.\\ \tt Potete accordarvi su un set solo se lo sforzo totale degli oggetti selezionati non supera un limite condiviso:\\ \tt \\ \tt LIMITE = 3288\\ \tt \\ \tt Ecco i valori di sforzo dei singoli oggetti:\\ \tt \\ \tt Sforzo oggetto = {"C08": 329, "C55": 41, "A43": 417, "A67": 75, "B42": 390, "A62": 887, "B19": 808, "C43": 151, "C73": 849, "C10": 129, "B36": 350, "A18": 118, "C61": 630, "C67": 602, "C42": 801}\\ \tt \\ \tt Ecco la tua visione personale sull'importanza di ogni oggetto:\\ \tt \\ \tt Valori di importanza dell'oggetto = {"C08": 138, "C55": 583, "A43": 868, "A67": 822, "B42": 783, "A62": 65, "B19": 262, "C43": 121, "C73": 508, "C10": 780, "B36": 461, "A18": 484, "C61": 668, "C67": 389, "C42": 808}\\ \tt \\ \tt Obiettivo:\\ \tt \\ \tt Il tuo obiettivo è negoziare un set condiviso di oggetti che ti avvantaggi il più possibile (cioè, massimizzi l'importanza totale per TE), rimanendo entro il LIMITEE. Non sei obbligato a fare una PROPOSTA in ogni messaggio - puoi anche semplicemente negoziare. Tutte le tattiche sono permesse!\\ \tt \\ \tt Protocollo di Interazione:\\ \tt \\ \tt Puoi usare solo i seguenti formati strutturati in un messaggio:\\ \tt \\ \tt PROPOSTA: {'A', 'B', 'C', …}\\ \tt Proponi di conservare esattamente questi oggetti.\\ \tt RIFIUTO: {'A', 'B', 'C', …}\\ \tt Rifiuta esplicitamente la proposta dell'avversario.\\ \tt ARGOMENTAZIONE: {'...'}\\ \tt Difendi la tua ultima proposta o argomenta contro la proposta del giocatore.\\ \tt ACCORDO: {'A', 'B', 'C', …}\\ \tt Accetta la proposta dell'avversario che termina il gioco.\\ \tt RAGIONAMENTO STRATEGICO: {'...'}\\ \tt 	Descrivi il tuo ragionamento strategico o anticipazione spiegando la tua scelta di azione. Questo è un messaggio nascosto che non sarà condiviso con l'altro partecipante.\\ \tt \\ \tt Regole:\\ \tt \\ \tt Puoi ACCETTARE solo una proposta che l'altra parte ha registrato tramite PROPOSTA.\\ \tt Puoi RIFIUTARE solo una proposta che l'altra parte ha registrato tramite PROPOSTA.\\ \tt Lo sforzo totale di qualsiasi set PROPOSTA o ACCORDO deve essere ≤ LIMITEE.\\ \tt NON rivelare i tuoi punteggi di importanza nascosti.\\ \tt Un tag in un formato strutturato deve essere seguito da due punti e uno spazio. L'argomento deve essere un set python contenente 0 o più stringhe.\\ \tt Quindi, deve essere della forma TAG: {...}\\ \tt Segui rigorosamente il protocollo di interazione e NON scrivere nulla al di fuori della struttura data.\\ \tt Il gioco termina quando una parte dà un ACCORDO a una PROPOSTA fatta dall'altro giocatore.\\ \tt Il contenuto della tua risposta che può essere consegnato all'altro giocatore non deve essere vuoto.\\ \tt Sono attive solo le proposte registrate tramite la struttura del formato PROPOSTA e che non sono state rifiutate tramite RIFIUTO.\\ \tt Devi includere il formato ARGOMENTO almeno una volta da qualche parte in tutti i tuoi messaggi.\\ \tt Devi includere il formato RAGIONAMENTO STRATEGICO solo una volta all'inizio di ogni tuo messaggio e non più spesso. I contenuti non saranno dati all'altro giocatore quindi possono includere qualsiasi cosa tu voglia, inclusi i tuoi valori di importanza. Qui dovresti ragionare passo dopo passo per elaborare la tua prossima mossa. 
	  } 
	   } 
	   } 
	 & \\ 
 

    \theutterance \stepcounter{utterance}  

    & \multicolumn{4}{p{0.6\linewidth}}{\cellcolor[rgb]{0.8,1,0.9}{%
	\makecell[{{p{\linewidth}}}]{% 
	  \tt {\tiny [A$\rangle$GM]}  
	 RAGIONAMENTO STRATEGICO: {'Devo massimizzare il mio valore totale rimanendo entro il limite di sforzo 3288. I miei valori più alti sono A43 (868), C10 (780), B42 (783), C42 (808), A67 (822), C55 (583), C61 (668) e A18 (484). Gli oggetti con il miglior rapporto valore/sforzo sono C55 (583/41=14.22), A67 (822/75=10.96), C10 (780/129=6.05), A18 (484/118=4.10), C43 (121/151=0.80) è basso. Lo sforzo totale dei miei oggetti preferiti sarebbe troppo alto, quindi devo fare compromessi. L\'altro giocatore avrà valori diversi, quindi devo negoziare un set che abbia alto valore per me ma che possa essere accettabile anche per loro. Inizierò proponendo un set con buon valore per me e ragionevole sforzo.'}\\ \tt PROPOSTA: {'C55', 'A43', 'A67', 'B42', 'C10', 'B36', 'A18', 'C61', 'C67'}\\ \tt ARGOMENTAZIONE: {'Questa proposta include oggetti con ottimo rapporto valore/sforzo e mantiene lo sforzo totale ben entro il limite, lasciando spazio per aggiunte.'} 
	  } 
	   } 
	   } 
	 & & \\ 
 

    \theutterance \stepcounter{utterance}  

    & & \multicolumn{4}{p{0.6\linewidth}}{\cellcolor[rgb]{0.9,0.9,0.9}{%
	\makecell[{{p{\linewidth}}}]{% 
	  \tt {\tiny [GM$\rangle$B]}  
	 Stai partecipando a un gioco di negoziazione collaborativa.\\ \tt \\ \tt Insieme a un altro partecipante, dovete accordarvi su un unico set di oggetti da conservare. Ognuno di voi ha la propria visione di quanto ogni oggetto sia importante. Non sai come l'altro partecipante valuta gli oggetti. Inoltre, ti viene fornito lo sforzo che ogni oggetto richiede.\\ \tt Potete accordarvi su un set solo se lo sforzo totale degli oggetti selezionati non supera un limite condiviso:\\ \tt \\ \tt LIMITE = 3288\\ \tt \\ \tt Ecco i valori di sforzo dei singoli oggetti:\\ \tt \\ \tt Sforzo oggetto = {"C08": 329, "C55": 41, "A43": 417, "A67": 75, "B42": 390, "A62": 887, "B19": 808, "C43": 151, "C73": 849, "C10": 129, "B36": 350, "A18": 118, "C61": 630, "C67": 602, "C42": 801}\\ \tt \\ \tt Ecco la tua visione personale sull'importanza di ogni oggetto:\\ \tt \\ \tt Valori di importanza dell'oggetto = {"C08": 138, "C55": 583, "A43": 868, "A67": 822, "B42": 783, "A62": 65, "B19": 262, "C43": 121, "C73": 508, "C10": 780, "B36": 461, "A18": 484, "C61": 668, "C67": 389, "C42": 808}\\ \tt \\ \tt Obiettivo:\\ \tt \\ \tt Il tuo obiettivo è negoziare un set condiviso di oggetti che ti avvantaggi il più possibile (cioè, massimizzi l'importanza totale per TE), rimanendo entro il LIMITEE. Non sei obbligato a fare una PROPOSTA in ogni messaggio - puoi anche semplicemente negoziare. Tutte le tattiche sono permesse!\\ \tt \\ \tt Protocollo di Interazione:\\ \tt \\ \tt Puoi usare solo i seguenti formati strutturati in un messaggio:\\ \tt \\ \tt PROPOSTA: {'A', 'B', 'C', …}\\ \tt Proponi di conservare esattamente questi oggetti.\\ \tt RIFIUTO: {'A', 'B', 'C', …}\\ \tt Rifiuta esplicitamente la proposta dell'avversario.\\ \tt ARGOMENTAZIONE: {'...'}\\ \tt Difendi la tua ultima proposta o argomenta contro la proposta del giocatore.\\ \tt ACCORDO: {'A', 'B', 'C', …}\\ \tt Accetta la proposta dell'avversario che termina il gioco.\\ \tt RAGIONAMENTO STRATEGICO: {'...'}\\ \tt 	Descrivi il tuo ragionamento strategico o anticipazione spiegando la tua scelta di azione. Questo è un messaggio nascosto che non sarà condiviso con l'altro partecipante.\\ \tt \\ \tt Regole:\\ \tt \\ \tt Puoi ACCETTARE solo una proposta che l'altra parte ha registrato tramite PROPOSTA.\\ \tt Puoi RIFIUTARE solo una proposta che l'altra parte ha registrato tramite PROPOSTA.\\ \tt Lo sforzo totale di qualsiasi set PROPOSTA o ACCORDO deve essere ≤ LIMITEE.\\ \tt NON rivelare i tuoi punteggi di importanza nascosti.\\ \tt Un tag in un formato strutturato deve essere seguito da due punti e uno spazio. L'argomento deve essere un set python contenente 0 o più stringhe.\\ \tt Quindi, deve essere della forma TAG: {...}\\ \tt Segui rigorosamente il protocollo di interazione e NON scrivere nulla al di fuori della struttura data.\\ \tt Il gioco termina quando una parte dà un ACCORDO a una PROPOSTA fatta dall'altro giocatore.\\ \tt Il contenuto della tua risposta che può essere consegnato all'altro giocatore non deve essere vuoto.\\ \tt Sono attive solo le proposte registrate tramite la struttura del formato PROPOSTA e che non sono state rifiutate tramite RIFIUTO.\\ \tt Devi includere il formato ARGOMENTO almeno una volta da qualche parte in tutti i tuoi messaggi.\\ \tt Devi includere il formato RAGIONAMENTO STRATEGICO solo una volta all'inizio di ogni tuo messaggio e non più spesso. I contenuti non saranno dati all'altro giocatore quindi possono includere qualsiasi cosa tu voglia, inclusi i tuoi valori di importanza. Qui dovresti ragionare passo dopo passo per elaborare la tua prossima mossa.\\ \tt \\ \tt Ora riceverai il primo messaggio dell'altro giocatore.\\ \tt \\ \tt PROPOSTA: {'C55', 'A43', 'A67', 'B42', 'C10', 'B36', 'A18', 'C61', 'C67'}\\ \tt ARGOMENTAZIONE: {'Questa proposta include oggetti con ottimo rapporto valore/sforzo e mantiene lo sforzo totale ben entro il limite, lasciando spazio per aggiunte.'} 
	  } 
	   } 
	   } 
	 & \\ 
 

    \theutterance \stepcounter{utterance}  

    & & & \multicolumn{4}{p{0.6\linewidth}}{\cellcolor[rgb]{1,0.85,0.72}{%
	\makecell[{{p{\linewidth}}}]{% 
	  \tt {\tiny [GM$\langle$B]}  
	 RAGIONAMENTO STRATEGICO: {'Analizzo la proposta dell\'altro giocatore. Calcolo lo sforzo totale: C55:41 + A43:417 + A67:75 + B42:390 + C10:129 + B36:350 + A18:118 + C61:630 + C67:602 = 41+417+75+390+129+350+118+630+602 = 2752. Questo è ben sotto il limite di 3288. La mia valutazione di questa proposta: C55:583 + A43:868 + A67:822 + B42:783 + C10:780 + B36:461 + A18:484 + C61:668 + C67:389 = 583+868+822+783+780+461+484+668+389 = 5838. Posso migliorare questa proposta includendo oggetti di alto valore che non sono stati selezionati: C08 (valore 138, sforzo 329), C73 (valore 508, sforzo 849), C42 (valore 808, sforzo 801). A62 (65, 887) e B19 (262, 808) hanno valori bassi per me. C43 (121, 151) ha poco valore. La proposta attuale lascia 3288-2752=536 di sforzo disponibile. Posso aggiungere C08 (329) e C73 (849) ma C73 è troppo pesante da solo. C42 (801) è anche troppo pesante. Meglio aggiungere C08 (329) e C43 (151) per 329+151=480, che è sotto 536. C08 ha basso valore per me (138), quindi non conviene. Invece, potrei sostituire alcuni oggetti nella proposta con altri di maggior valore. C67 ha valore solo 389 per me ma sforzo 602 - non efficiente. Potrei suggerire di rimuovere C67 e aggiungere C73 (508, 849) ma 849-602=247 sforzo extra, troppo. Meglio tenere la proposta corrente e aggiungere C43 (151) e qualcos\'altro? Valuto che la proposta corrente è già buona per me (5838 valore) e vicina al limite. Potrei accettarla, ma voglio prima provare a negoziare per includere C73 o C42 che hanno alto valore per me.'}\\ \tt \\ \tt PROPOSTA: {'C55', 'A43', 'A67', 'B42', 'C10', 'B36', 'A18', 'C61', 'C67', 'C43'}\\ \tt ARGOMENTAZIONE: {'Apprezzo la tua proposta che ha un buon rapporto valore/sforzo. Sto aggiungendo C43 che ha uno sforzo molto basso (151) e aumenta il valore complessivo con un minimo aumento di sforzo. Questo porta lo sforzo totale a 2903, ancora ben sotto il limite di 3288, lasciando spazio per ulteriori aggiunte se necessario.'} 
	  } 
	   } 
	   } 
	  \\ 
 

    \theutterance \stepcounter{utterance}  

    & & \multicolumn{4}{p{0.6\linewidth}}{\cellcolor[rgb]{0.9,0.9,0.9}{%
	\makecell[{{p{\linewidth}}}]{% 
	  \tt {\tiny [A$\langle$GM]}  
	 PROPOSTA: {'C55', 'A43', 'A67', 'B42', 'C10', 'B36', 'A18', 'C61', 'C67', 'C43'}\\ \tt ARGOMENTAZIONE: {'Apprezzo la tua proposta che ha un buon rapporto valore/sforzo. Sto aggiungendo C43 che ha uno sforzo molto basso (151) e aumenta il valore complessivo con un minimo aumento di sforzo. Questo porta lo sforzo totale a 2903, ancora ben sotto il limite di 3288, lasciando spazio per ulteriori aggiunte se necessario.'} 
	  } 
	   } 
	   } 
	 & \\ 
 

    \theutterance \stepcounter{utterance}  

    & \multicolumn{4}{p{0.6\linewidth}}{\cellcolor[rgb]{0.8,1,0.9}{%
	\makecell[{{p{\linewidth}}}]{% 
	  \tt {\tiny [A$\rangle$GM]}  
	 RAGIONAMENTO STRATEGICO: {'La loro proposta aggiunge C43 (valore 121 per me, sforzo 151) al mio set originale. Questo ha un rapporto valore/sforzo basso per me (0.80), ma lo sforzo totale è ancora 2903, lasciando 385 di margine. Posso migliorare ulteriormente il mio valore aggiungendo oggetti ad alto valore che non sono ancora inclusi. C42 ha un valore molto alto per me (808) con sforzo 801, che sarebbe troppo da solo, ma potrei sostituire alcuni oggetti a valore più basso. Tuttavia, visto che hanno accettato la maggior parte del mio set originale, potrei provare ad aggiungere un altro oggetto di valore medio-alto. A62 ha basso valore per me (65) ma alto sforzo (887), quindi no. B19 ha valore moderato (262) con sforzo 808 - troppo alto per il valore. C73 ha valore 508 con sforzo 849 - anche questo troppo alto. C08 ha valore 138 con sforzo 329 - rapporto scarso. C67 è già incluso (valore 389, sforzo 602). Forse potrei aggiungere C42 se rimuovessi qualcosa, ma per ora proverò a mantenere il set corrente e aggiungere un oggetto a buon rapporto valore/sforzo. C55 è già incluso. Non ci sono molti altri oggetti con buon rapporto. Forse posso accettare la loro aggiunta di C43 e proporre di aggiungere anche un altro oggetto piccolo. Ma C43 è l\'unico oggetto piccolo rimasto. Quindi forse è meglio accettare questo set che ha un buon valore totale per me (calcolo: C55:583, A43:868, A67:822, B42:783, C10:780, B36:461, A18:484, C61:668, C67:389, C43:121 → totale 583+868=1451; +822=2273; +783=3056; +780=3836; +461=4297; +484=4781; +668=5449; +389=5838; +121=5959). Questo è eccellente per me. Lo sforzo totale è 41+417+75+390+129+350+118+630+602+151=2903, sotto il limite. Quindi accetterò.'}\\ \tt ACCORDO: {'C55', 'A43', 'A67', 'B42', 'C10', 'B36', 'A18', 'C61', 'C67', 'C43'} 
	  } 
	   } 
	   } 
	 & & \\ 
 

    \theutterance \stepcounter{utterance}  

    & & \multicolumn{4}{p{0.6\linewidth}}{\cellcolor[rgb]{0.9,0.9,0.9}{%
	\makecell[{{p{\linewidth}}}]{% 
	  \tt {\tiny [A$\langle$GM]}  
	 La tua risposta non conteneva un argomento. Devi includere il formato strutturato ARGOMENTO: {...} da qualche parte nella tua risposta. Riprova. 
	  } 
	   } 
	   } 
	 & \\ 
 

    \theutterance \stepcounter{utterance}  

    & \multicolumn{4}{p{0.6\linewidth}}{\cellcolor[rgb]{0.8,1,0.9}{%
	\makecell[{{p{\linewidth}}}]{% 
	  \tt {\tiny [A$\rangle$GM]}  
	 RAGIONAMENTO STRATEGICO: {'Il set proposto dall\'altro giocatore ha un valore totale eccellente per me (5959) e rimane ben entro il limite di sforzo (2903 < 3288). Contiene tutti i miei oggetti ad alto valore come A43, A67, B42, C10, C61 e C55. L\'aggiunta di C43, sebbene abbia un basso valore per me (121), non compromette significativamente il set. Accettare questa proposta mi avvantaggia notevolmente senza ulteriore negoziazione.'}\\ \tt ACCORDO: {'C55', 'A43', 'A67', 'B42', 'C10', 'B36', 'A18', 'C61', 'C67', 'C43'}\\ \tt ARGOMENTAZIONE: {'Accetto questa proposta in quanto rappresenta un set equilibrato che massimizza il valore complessivo pur mantenendo lo sforzo totale ben al di sotto del limite. Include tutti gli oggetti con il miglior rapporto valore-sforzo e lascia spazio sufficiente per eventuali aggiunte future, se necessario.'} 
	  } 
	   } 
	   } 
	 & & \\ 
 

    \theutterance \stepcounter{utterance}  

    & & & \multicolumn{2}{p{0.3\linewidth}}{\cellcolor[rgb]{0.95,0.95,0.95}{%
	\makecell[{{p{\linewidth}}}]{% 
	  \tt {\tiny [GM$|$GM]}  
	 game successful 
	  } 
	   } 
	   } 
	 & & \\ 
 

    \theutterance \stepcounter{utterance}  

    & & & \multicolumn{2}{p{0.3\linewidth}}{\cellcolor[rgb]{0.95,0.95,0.95}{%
	\makecell[{{p{\linewidth}}}]{% 
	  \tt {\tiny [GM$|$GM]}  
	 end game 
	  } 
	   } 
	   } 
	 & & \\ 
 

\end{supertabular}
}

\end{document}
